% Zdroj: PDF priklady-jaro-2026.pdf, fyzická str. 3. Značka: společné okruhy. Zařazeno: M2.
\section{Matice}
\szzotazka{jaro 2026}{3}{Matice}

\noindent\textbf{Zadání.}\par
Uvažujme matici $A \in \mathbb{R}^{3 \times 3}$, jejíž řádkový prostor má bázi tvořenou pouze vektorem $(1,2,3)^T$.
\begin{enumerate}
  \item Najděte odstupňovaný tvar matice $A$.
  \item Napište definici determinantu a spočítejte konkrétní hodnotu $\det(A)$.
  \item Najděte ortogonální bázi jádra matice $A$.
\end{enumerate}

\medskip
\noindent\textbf{Náčrt řešení.}\par
\begin{enumerate}
  \item
    \[
      \begin{pmatrix}
        1 & 2 & 3 \\
        0 & 0 & 0 \\
        0 & 0 & 0
      \end{pmatrix}.
    \]
  \item
    \[
      \det(A) = \sum_{p \in S_n} \operatorname{sgn}(p) \prod_{i=1}^{n} a_{i,p(i)}.
    \]
    V našem případě má matice $A$ hodnost 1, čili je singulární a $\det(A) = 0$.
  \item Jádro matice $A$ má dimenzi 2 a bázi tvoří např.\ vektory $(2,-1,0)^T$, $(3,0,-1)^T$. Gramovou--Schmidtovou ortogonalizací dostaneme ortogonální bázi $(2,-1,0)^T$, $(3,6,-5)^T$.
\end{enumerate}

\medskip
\noindent\textit{Doplňující vysvětlení (nad rámec oficiálního náčrtu):} (1)~Všechny řádky $A$ jsou násobky $(1,2,3)$ a~aspoň jeden je nenulový; eliminací (a~vydělením pivotu) zbude jediný nenulový řádek $(1,2,3)$. (3)~Jádro je množina $x$ s~$x_1+2x_2+3x_3=0$ (dimenze $3-1=2$). Gram--Schmidt na $v_1=(2,-1,0)^T$, $v_2=(3,0,-1)^T$: $\langle v_2,v_1\rangle=6$, $\langle v_1,v_1\rangle=5$, takže $v_2-\tfrac65v_1=\tfrac15(3,6,-5)^T$; po vynásobení pěti dostaneme $(3,6,-5)^T$. Zkouška: $(2,-1,0)\cdot(3,6,-5)=0$ a~$3+2\cdot6+3\cdot(-5)=0$, vektor tedy leží v~jádře.
